% Standalone article section; uses amsmath, amssymb and amsthm environments.
% Citation keys are supplied in literature-references.bib.
\section{Reconstructing an unmarked cut manifold}
\label{sec:unmarked-reconstruction}

Large normal coordinates do not by themselves force a large triangulation of
the manifold obtained by cutting.  A long stack of parallel normal discs is
represented by an interval bundle, whose topological complexity may be much
smaller than the number of discs in the stack.  There are nevertheless two
distinct costs: the topology of the interval-bundle bases, and the information
in the attaching maps.  The purpose of this section is to state an explicit
reconstruction bound that accounts for both.

The result below is a constructive consequence of surface classification and
the published parallelity-bundle machinery.  We do not claim it as a new major
topological theorem.  Its contribution here is to make the output convention,
the cellular input requirement, and the limits of its use in a hierarchy
algorithm explicit.

For any compact surface $F$, orientable or otherwise, set
\[
  \chi_-(F)=\sum_{F_0\in\pi_0(F)}\max\{0,-\chi(F_0)\}.
\]
The word \emph{unmarked} means that the output is required to have the correct
homeomorphism type, without a prescribed parametrisation of its boundary or
a prescribed identification with a subspace of the original manifold.

\subsection{A reconstruction bound with a supplied cellular core}

\begin{theorem}[Reconstruction from explicit core and bundle data]
\label{thm:cellular-core-reconstruction}
Suppose that a compact orientable $3$-manifold, with corners allowed, is
specified as
\[
  X=C\mathbin{\cup}\bigcup_{j=1}^{m}B_j,
  \qquad B_j\longrightarrow F_j,
\]
where the $F_j$ are compact connected surfaces and the $B_j$ are pairwise
disjoint interval bundles with orientable total spaces.  Assume the following
input data.
\begin{enumerate}
  \item The core $C$ has an explicit cellular decomposition with total size
  at most $q$, using a fixed finite collection of $3$-ball templates.  The
  incidence data and the cellular attaching maps are included in this size.
  Equivalently, one may supply an explicit triangulation of $C$ with $O(q)$
  tetrahedra, together with its boundary subdivisions.
  \item Each base $F_j$ is specified by its orientability, genus or crosscap
  number, and number of boundary components.  The interval-bundle type is
  specified.  There are no identifications between horizontal boundaries of
  the $B_j$ in this presentation.
  \item Each $B_j\cap C$ is a subset of the vertical boundary
  $\partial_vB_j$.  Its attaching data consist of explicitly ordered fibre
  rectangles and their cellular identifications with rectangles on
  $\partial C$.  Boundary gaps, horizontal-sheet incidence, and the relevant
  orientations are recorded.  The total size of this data and of the
  required subdivisions of $\partial C$ is included in $q$.
\end{enumerate}
The presentation is assumed to be valid; recognising validity of arbitrary
cellular presentations is not part of the assertion.  Then $X$ admits an
ordinary, unmarked triangulation with
\[
  O\!\left(q+\sum_{j=1}^{m}(1+\chi_-(F_j))\right)
  \tag{\ref{thm:cellular-core-reconstruction}.1}
\]
tetrahedra.  Such a triangulation can be constructed in time polynomial in
the input bit length and the displayed bound.

Every component and every repeated copy is counted in the sum.  If a family
of detached, identical bundles is instead specified by a binary
multiplicity, expanding that family incurs the corresponding output cost.
\end{theorem}

\begin{proof}
First refine the fixed core templates and their attaching rectangles into
triangles and tetrahedra.  A fixed template requires a fixed number of
tetrahedra; the stipulated cellular size includes all necessary boundary
subdivisions.  Consequently this stage has cost $O(q)$.

We use the following elementary surface-construction fact.  A connected
compact surface $F$, together with $v$ prescribed subdivision vertices on
its boundary, has a triangulation respecting those vertices with
\[
  O(1+\chi_-(F)+v)
\]
triangles.  To see this, choose the usual polygon model from the classification
of compact surfaces, triangulate it, and subdivide the boundary as prescribed.
For an orientable surface with $b$ boundary components,
$\chi(F)=2-2g-b$; in the nonorientable case,
$\chi(F)=2-k-b$, where $k$ is the crosscap number.  These formulas bound the
size of a polygon model by $O(1+\chi_-(F))$, apart from the prescribed boundary
subdivisions.  The sphere, disc, annulus, projective plane, and Klein bottle
are covered by fixed templates.  A bounded subdivision converts any polygon
identifications into a genuine surface triangulation.  Prescribed cyclic
orders on boundary components are respected throughout.

Apply this construction to each $F_j$, introducing the subdivision points
dictated by the ordered attaching rectangles.  Their total number is $O(q)$.
Boundary components with no attaching rectangles require only a bounded
number of additional vertices each; their number is already bounded by
$O(1+\chi_-(F_j))$.  Thus the total number of base triangles is
\[
  O\!\left(q+\sum_j(1+\chi_-(F_j))\right).
\]

An interval bundle over a triangle is a triangular prism.  Glue these prisms
using the specified interval-bundle type.  For an orientable total space,
the twisting character of the bundle is the orientation character of its
base: the bundle is a product over an orientable base, and is the orientable
twisted interval bundle over a nonorientable base.  The polygon construction
therefore supplies the needed transition maps directly.  Subdivide prism
faces compatibly and triangulate every prism with a bounded number of
tetrahedra.  For example, first barycentrically subdivide the polygonal faces
and then cone each subdivided prism boundary to an interior vertex; this
avoids choices of incompatible diagonals.  The attaching rectangles match
the refined core cells after a bounded further subdivision.  Gluing gives
the required triangulation of $X$, with the asserted bound.

All constructions are finite operations on explicit cells, together with
integer arithmetic on the supplied topological parameters.  Writing down
the polygon models and the tetrahedra takes time polynomial in the input
bit length plus the output size.  Binary multiplicities can be retained in
a compressed representation; if the requested output is an ordinary
triangulation, the corresponding copies must be written out.
\end{proof}

\paragraph{Why counting handles alone is insufficient.}
The hypothesis is a bound on \emph{cellular size}, including attaching
maps.  It is not a bound on the number of abstract handles whose attaching
words might have unbounded length.  An arbitrary gluing map given by a short
word in large powers of mapping classes does not satisfy the hypothesis
merely because its word has few syllables.  The theorem applies to the
explicit ordered fibre rectangles in its statement.

There is no additional unrecorded torus-slope parameter in this particular
presentation: the bundles attach through their vertical annuli and
rectangles, while their horizontal boundaries remain unglued.  For example,
a Dehn twist on a vertical annulus that fixes its boundary extends over a
collar of that annulus in the bundle.  In collar coordinates
$S^1\times[0,\varepsilon]\times[0,1]$, with the vertical annulus at $r=0$,
an extension of $k$ twists is
\[
  (\theta,r,s)\longmapsto
  (\theta+2\pi k s f(r),r,s),
  \qquad f(0)=1,\quad f(\varepsilon)=0.
\]
It changes the parametrisation on a horizontal boundary collar.  This is
permitted for an unmarked output, and explains exactly why the same argument
does not provide a reconstruction relative to fixed horizontal seam charts.

\subsection{Controlling the bases by the cutting surface}

The following elementary estimate converts the topological size of the
bundle bases into a parameter of the cutting surface.

\begin{lemma}[Subsurface complexity]
\label{lem:subsurface-complexity}
Let $\Sigma$ be a compact surface and let $A\subset\Sigma$ be a compact
subsurface of dimension two.  After rounding corners, suppose the relative frontier of $A$
in $\Sigma$ is a disjoint union of $c$ circles and properly embedded arcs.
Write
\[
  \kappa(Y)=b_0(Y)+b_1(Y;\mathbb Q)
\]
for any compact surface $Y$.  Then
\[
  b_0(A)\leq b_0(\Sigma)+c,
  \qquad
  b_1(A;\mathbb Q)\leq b_1(\Sigma;\mathbb Q)+c,
\]
and consequently
\[
  \kappa(A)\leq\kappa(\Sigma)+2c.
\]
If $\Sigma$ is orientable, then also
\[
  \kappa(\Sigma)\leq\chi_-(\Sigma)+3b_0(\Sigma).
\]
\end{lemma}

\begin{proof}
Cutting a surface along one properly embedded arc or circle increases the
number of components by at most one.  The closures of the components on the
$A$ side form a subcollection of the components obtained by cutting along
the frontier.  This proves the bound for $b_0(A)$.

Let $D$ be the closure of the complement of $A$.  Use small collar
neighbourhoods to obtain a Mayer--Vietoris decomposition of $\Sigma$ whose
intersection retracts onto the frontier.  Its first rational Betti number
is at most $c$.  Exactness of
\[
 H_1(A\cap D;\mathbb Q)\longrightarrow
 H_1(A;\mathbb Q)\oplus H_1(D;\mathbb Q)
 \longrightarrow H_1(\Sigma;\mathbb Q)
\]
shows that
$b_1(A;\mathbb Q)+b_1(D;\mathbb Q)
\leq b_1(\Sigma;\mathbb Q)+c$.
Discard the nonnegative term $b_1(D;\mathbb Q)$ and add the component bound.
Finally, the classification formulas for a connected orientable surface,
including the closed and nonclosed cases, give
$1+b_1\leq\chi_-+3$.  Summation over components completes the proof.
\end{proof}

\begin{corollary}[Size of an unmarked normal cut]
\label{cor:normal-cut-reconstruction}
Let $M$ be a compact orientable $3$-manifold given by a triangulation with
$t$ tetrahedra.  Let $S$ be a compact orientable normal surface of weight
$W$, specified by its binary normal coordinates.  Retain all components of
\[
  M_S=M\setminus\operatorname{int}N(S).
\]
Then an ordinary unmarked triangulation of $M_S$ with
\[
  O\bigl(t+b_0(S)+\chi_-(S)\bigr)
\]
tetrahedra can be constructed in time polynomial in
\[
  t,\quad \log(2+W),\quad b_0(S)+\chi_-(S).
\]
No irreducibility or essentialness assumption is needed for this cutting
statement.  The running-time assertion includes the cost of expanding all
output components.
\end{corollary}

\begin{proof}
Use the parallelity bundle $\mathcal B$ for the pair consisting of $M_S$
and the two copies $S^+\sqcup S^-$ of $S$ in its boundary.  The compressed-cut
algorithm of Lackenby gives the remaining core, the base-surface types, and
ordered vertical attachment data in time polynomial in
$t\log(2+W)$; see Theorem~9.3 and the explanatory paragraphs and proof in
Section~9.4 of \cite{Lackenby2021Certification}.  In a tetrahedral source,
the nonparallel pieces and attachment squares have bounded complexity in
each local tetrahedral cell.  Summing over these cells gives $O(t)$ total
cellular incidences, including the subdivisions of vertical attachments.
An individual dual two-handle may have many incidences when its original
edge has high valence; only the total incidence bound is asserted here.
Accordingly, the construction supplies the explicit core and attachment
complexity $q=O(t)$ needed in
Theorem~\ref{thm:cellular-core-reconstruction}; no bound on handle count alone
is being substituted.  Repeated detached components are retained with
multiplicities until expansion is requested.

Set $A=\partial_h\mathcal B\subset S^+\sqcup S^-$ and let $F_j$ be the
connected base surfaces.  The relative frontier of $A$ in
$S^+\sqcup S^-$ is incident to the explicit nonparallel core.  Each of its
arcs or circles contains an edge of the vertical attachment data, and there
are only $O(t)$ such edges.  Portions lying in the original boundary of
$S^+\sqcup S^-$ are not part of this relative frontier.  Thus its number
$c$ of components is $O(t)$.  Components of $\mathcal B$ with no core
attachment are allowed in this argument: they contribute whole components
to $A$, rather than additional frontier.

Apply Lemma~\ref{lem:subsurface-complexity} to obtain
\[
  \kappa(A)
  \leq 2\kappa(S)+O(t)
  =O\bigl(t+b_0(S)+\chi_-(S)\bigr).
\]
The horizontal boundary $A_j$ of the bundle over $F_j$ is its two-sheeted
boundary cover.  Hence
\[
  \chi_-(A_j)=2\chi_-(F_j),
  \qquad b_0(A_j)\geq 1.
\]
For every compact surface, $\chi_-$ is at most its first rational Betti
number.  Consequently
\[
  \sum_j(1+\chi_-(F_j))
  \leq b_0(A)+\tfrac12\chi_-(A)
  \leq\kappa(A).
\]
Theorem~\ref{thm:cellular-core-reconstruction} now gives the tetrahedron
bound and the construction time.  The weighted orbit machinery can keep
identical component data grouped with binary multiplicities throughout the
compressed part of the computation; writing out their copies costs only
the explicit output parameter in the statement.
\end{proof}

\paragraph{A necessary component term.}
The term $b_0(S)$ cannot simply be dropped from a theorem retaining the whole
cut manifold.  Fix a triangulation of a $3$-ball having an interior vertex,
and let $v$ be the normal vector of its vertex-linking sphere.  The vector
$Nv$ represents $N$ nested parallel normal spheres and has bit length
$O(\log(2+N))$.  Cutting along them produces $N+1$ nonempty $3$-manifold
components.  An ordinary triangulation of their disjoint union contains at
least $N+1$ tetrahedra, although every cutting component has
$\chi_-=0$.  One may retain these product regions by a compressed
multiplicity, or expand them and pay their output size; these are different
output specifications.

\paragraph{What this does and does not supply to a hierarchy.}
For a connected disc, annulus, sphere, or torus, the displayed unmarked size
bound is $O(t)$.  More generally, it replaces dependence on raw normal
weight by dependence on the topological size of the cutting surface.  This
is useful when a prospective hierarchy supplies a bound on
$b_0(S)+\chi_-(S)$ even though the normal coordinates are very large.
Repeatedly applying the corollary, however, does not by itself give a
hierarchy algorithm of quasipolynomial complexity.  It supplies neither a
bound on the number of cuts nor a controlled identification of their
horizontal boundaries with previously stored surfaces.  In particular,
transferring an existing boundary pattern, preserving normal-coordinate
charts for later cuts, and retaining the pairing that reglues
$S^+$ to $S^-$ require additional information.  A future implementation
should therefore expose two separate interfaces: unmarked reconstruction
with the bound proved here, and reconstruction relative to explicitly
costed boundary markings.  The latter remains a research obligation.
